% Clustering -- parent section; each part is a subsection in its own file.

\section{Clustering}                     % phenotype discovery in the latent space
\label{sec:clustering}                   % \secref{sec:clustering} -- used by Encoding and Morphological Metrics

The implicit model of represents every airway tree by a latent code. If
groups of trees with a similar shape exist, they should appear as groups of nearby codes. Following
the phenotype-discovery approach of \citet{norouzi2024volpam}, who cluster the latent codes of a
3D autoencoder to find imaging phenotypes without any labels, we cluster the latent codes of the
airway trees with two families of methods: agglomerative hierarchical clustering and K-means,
the latter also in a kernel version. The groups found are then visualized with three
projections of the latent space and characterized with the morphological descriptors of
\secref{sec:morphometrics} (Figure~\ref{fig:clustering_pipeline}).

\begin{figure}[htbp]                     % the clustering pipeline at a glance
  \centering                             % center the diagram
  \begin{tikzpicture}[
      node distance=5mm and 5mm,                         % vertical and horizontal gaps
      box/.style={draw, rounded corners=2pt, align=center, text width=3.5cm,
                  inner sep=4pt, font={\small\hyphenpenalty=10000}},  % one box per step, no hyphens
      arr/.style={-{Stealth[length=2mm]}}]               % arrow style
    \node[box] (codes) {Latent codes\\$\boldsymbol{\alpha}_i \in \mathbb{R}^{256}$, 419 trees};
    \node[box, right=of codes] (clus) {Clustering\\Agglomerative\\K-means, kernel K-means};
    \node[box, right=of clus, text width=2.5cm] (lab) {Cluster label\\for every tree};
    \node[box, above right=2mm and 5mm of lab] (proj) {Projections\\kernel PCA, t-SNE,\\UMAP};
    \node[box, below right=2mm and 5mm of lab] (morph) {Characterization\\morphological descriptors};
    \draw[arr] (codes) -- (clus);                         % codes are clustered
    \draw[arr] (clus) -- (lab);                           % clustering gives labels
    \draw[arr] (lab.east) -- (proj.west);                 % labels coloured in the projections
    \draw[arr] (lab.east) -- (morph.west);                % labels compared with the descriptors
  \end{tikzpicture}
  \caption{Clustering pipeline. The latent codes are clustered in their full space; the
           projections are only used to look at the result, and the morphological descriptors
           to describe what distinguishes the clusters.}
  \label{fig:clustering_pipeline}        % referenced as Figure~\ref{fig:clustering_pipeline}
\end{figure}

% % % Input of the clustering: the DGCI latent codes. Subsection of "Clustering".

% \subsection{Input: the latent codes}     % what is clustered
% \label{sec:clust_input}                  % \secref{sec:clust_input}

% After training, every airway tree $i$ is represented by its instance latent code
% $\boldsymbol{\alpha}_i \in \mathbb{R}^{128}$ (\secref{sec:dgci_training}). Stacking the codes of
% the $N = 419$ trees gives the data matrix $X \in \mathbb{R}^{419 \times 128}$, whose rows
% $\mathbf{x}_i = \boldsymbol{\alpha}_i$ are the points to cluster. All distances below are
% Euclidean distances between these codes.

% \paragraph{Active dimensions}            % not every latent dimension carries information
% Not every latent dimension is used by the model: the regularization of the latent codes
% (Eq.~\eqref{eq:dgci_reg}) pulls them towards the shared code, so dimensions that do not help the
% reconstruction stay almost constant across trees. To measure how much each dimension varies, we
% use its median absolute deviation (MAD), a spread measure that is not driven by a few outlying
% shapes,
% \begin{equation}
%   \mathrm{MAD}_j = \operatorname{median}_i \big|\, x_{ij} - \operatorname{median}_k x_{kj} \,\big|,
%   \qquad j = 1, \dots, 128.
%   \label{eq:mad}                         % robust spread of latent dimension j
% \end{equation}
% Sorted on a logarithmic scale, the MAD values show a clear gap between active and nearly
% constant dimensions; with a threshold of $3 \times 10^{-4}$ placed in this gap, about 45
% dimensions are active. \TODO{State which input was used for each method: all 128 dimensions,
% or only the active ones. In the notebooks, agglomerative clustering uses all 128; in the
% K-means notebook the final assignment \texttt{codes\_active = enc\_output} also selects all 128.}
   % Subsection "Input: the latent codes"
% Agglomerative hierarchical clustering (Ward). Subsection of "Clustering".

\subsection{Agglomerative hierarchical clustering}  % bottom-up merging, Ward linkage
\label{sec:clust_agglo}                  % \secref{sec:clust_agglo}

% \paragraph{Principle}                    % bottom-up construction of a hierarchy
Agglomerative clustering builds a hierarchy of nested groups from the bottom up
\cite{norouzi2024volpam}. At the start, every tree is a cluster of its own. At each step, the two
closest clusters are merged, and the procedure repeats until a single cluster contains all the
trees. The sequence of merges forms a binary tree, the \emph{dendrogram}, in which the height of
each junction shows how far apart the two merged clusters were. Cutting the dendrogram with a
horizontal line at some height gives a partition into $K$ clusters, so one run of the algorithm
yields the clustering for every $K$ and shows how small groups nest inside larger ones. This is
what makes hierarchical methods attractive for discovering intrinsic groupings whose number is
not known in advance \cite{norouzi2024volpam}.

\paragraph{Linkage}                      % how "closest" is defined: Ward
What ``closest'' means is set by the linkage. \citet{norouzi2024volpam} merge the two
clusters whose centroids are closest. We use Ward's linkage \cite{ward1963}, which merges the
pair of clusters whose union increases the total within-cluster sum of squares the least. For
two clusters $A$ and $B$ with centroids $\boldsymbol{\mu}_A$ and $\boldsymbol{\mu}_B$, this
increase is
\begin{equation}
  \Delta(A, B) = \frac{|A|\,|B|}{|A| + |B|}\,
                 \big\|\boldsymbol{\mu}_A - \boldsymbol{\mu}_B\big\|_2^2 .
  \label{eq:ward}                        % cost of merging A and B
\end{equation}
Like the centroid linkage, Ward's criterion compares cluster centres, but the factor in front of
the distance penalizes merging two large clusters, so it favours compact clusters of similar
size. Since it greedily minimizes the same within-cluster sum of squares as K-means
(Eq.~\eqref{eq:kmeans}), the two methods can be compared directly.

\paragraph{Choosing the number of clusters}  % elbow on the average intra-cluster distance
To choose where to cut the dendrogram, we follow the elbow criterion of
\citet{norouzi2024volpam}. For every number of clusters $k = 1, \dots, 19$, we cut the hierarchy
into $k$ clusters and compute the average intra-cluster distance, the mean distance between two
trees of the same cluster, averaged over the clusters:
\begin{equation}
  D_{\mathrm{avg}}(k) = \frac{1}{k} \sum_{l=1}^{k}
     \frac{1}{N_l^2} \sum_{u=1}^{N_l} \sum_{v=1}^{N_l}
     \big\|\mathbf{x}_{ul} - \mathbf{x}_{vl}\big\|_2 ,
  \label{eq:davg}                        % mean pairwise distance inside each cluster
\end{equation}
where $N_l$ is the number of trees in cluster $l$ and $\mathbf{x}_{ul}$ is the code of its
$u$-th tree. Splitting a cluster always makes its parts more compact, so $D_{\mathrm{avg}}$ tends
to decrease with $k$. The useful information is where it drops the most: the number of
clusters is chosen as the one with the largest relative reduction \cite{norouzi2024volpam},
\begin{equation}
  k^{*} = \arg\max_{k}\;
          \frac{D_{\mathrm{avg}}(k-1) - D_{\mathrm{avg}}(k)}{D_{\mathrm{avg}}(k-1)} .
  \label{eq:kstar}                       % largest relative drop of D_avg
\end{equation}
A large relative drop means that going from $k-1$ to $k$ clusters separates groups that are
genuinely different; after that point, further splits mostly cut coherent groups into pieces.
The elbow criterion identified $k^{\*}=4$ as the optimal number of clusters for both the registered and unregistered airway trees. Thus, in both representations, the hierarchy was cut at four clusters for the subsequent analysis.  % Subsection "Agglomerative hierarchical clustering"
% K-means. Subsection of "Clustering".

\subsection{K-means}
\label{sec:clust_kmeans}

Unlike agglomerative clustering, K-means directly partitions the trees into a fixed number $K$ of
clusters. Each cluster $c$ is represented by its centroid $\boldsymbol{\mu}_c$, and the
partition is chosen to minimize the \emph{inertia}, the total squared distance of the trees to
the centroid of their cluster:
\begin{equation}
  J = \sum_{i=1}^{N} \big\|\mathbf{x}_i - \boldsymbol{\mu}_{c(i)}\big\|_2^2 ,
  \label{eq:kmeans}
\end{equation}
where $c(i)$ is the cluster of tree $i$. Lloyd's algorithm \cite{lloyd1982} alternates two
steps: every tree is assigned to its nearest centroid, then every centroid is moved to the mean
of the trees assigned to it. Neither step can increase $J$, so the algorithm converges, but only
to a local minimum that depends on the initial centroids. We therefore initialize the centroids
with k-means++ \cite{arthur2007kmeanspp}, which spreads them out by picking distant trees with
higher probability, and keep the best of 20 runs with different initializations (scikit-learn
implementation \cite{pedregosa2011sklearn}). Because each tree joins its nearest centroid, the
borders between clusters are flat hyperplanes: K-means works best for compact, roughly
spherical clusters.

\paragraph{Choosing the number of clusters}
K-means requires the number of clusters $K$ to be specified in advance. We therefore evaluated
$K=2,\dots,10$ for each airway-tree representation. For each candidate $K$, K-means was fitted
to the 256-dimensional latent codes using k-means++ initialization with 20 restarts, retaining
the solution with the lowest inertia. Two criteria were then evaluated: the inertia, which
measures the within-cluster sum of squared distances, and the silhouette coefficient
\cite{rousseeuw1987silhouettes}, which measures how well each tree is separated from the other
clusters. The mean silhouette was used to assess the separation of the resulting partition.
\begin{equation}
  s(i) = \frac{b(i) - a(i)}{\max\{a(i),\, b(i)\}} \in [-1, 1].
  \label{eq:silhouette}
\end{equation}
The analysis identified $K=3$ as the preferred number of clusters for both the registered and
unregistered airway-tree representations. The final K-means analysis was therefore performed
with three clusters for both representations.         % Subsection "K-means and kernel K-means"
% Projections used to look at the latent space. Subsection of "Clustering".

\subsection{Visualizing the latent space}  % kernel PCA, t-SNE, UMAP
\label{sec:clust_projections}            % \secref{sec:clust_projections}

The latent codes are represented in a 256-dimensional space and therefore cannot be directly
visualized. We project them to two and three dimensions using three complementary dimensionality-
reduction methods and colour the resulting embeddings according to cluster membership or
morphological descriptors.
% The projections are used only for viewing: all clusters are computed in the full latent space and no projection preserves all
% distances, so a separation that appears in a projection, or the lack of one, has to be read with
% care. 
The three methods distort the data in different ways, which is why we compare them.

\paragraph{Kernel PCA}                   % linear PCA in the RBF feature space
Principal component analysis (PCA) finds the orthogonal directions of largest variance; kernel PCA
\cite{scholkopf1998kpca} does the same in the feature space of a kernel. With the RBF kernel of
Eq.~\eqref{eq:rbf}
\begin{equation}
  k(\mathbf{x}_i,\mathbf{x}_j) = \exp\!\left(-\gamma \|\mathbf{x}_i-\mathbf{x}_j\|^2\right)
  \label{eq:rbf}
\end{equation}
and the same median-rule $\gamma$ as kernel K-means, it shows the space in
which kernel K-means operates. The kernel matrix $K$ is first centred,
$\tilde{K} = H K H$ with $H = I - \frac{1}{N}\mathbf{1}\mathbf{1}^{\top}$, and decomposed as
$\tilde{K}\mathbf{v}_m = \lambda_m \mathbf{v}_m$. The coordinate of tree $i$ on the $m$-th
component is $\sqrt{\lambda_m}\, v_{m,i}$, and the share of kernel variance kept by a view with
$d$ components is $\sum_{m \le d} \lambda_m / \sum_m \lambda_m$. 
% Kernel PCA is deterministic and
% keeps the global layout of the data, but it only shows the few directions of largest variance.
% % \TODO{Report the kernel variance shown in 2D and in 3D (printed by the notebook).}

\paragraph{t-SNE}                        % neighbourhood probabilities, KL divergence
t-distributed stochastic neighbour embedding (t-SNE) \cite{vandermaaten2008tsne} preserves
neighbourhoods rather than distances. For every tree $i$, it turns distances into probabilities
of picking tree $j$ as a neighbour,
\begin{equation}
  p_{j|i} = \frac{\exp\!\big(-\|\mathbf{x}_i - \mathbf{x}_j\|^2 / 2\sigma_i^2\big)}
                 {\sum_{k \neq i} \exp\!\big(-\|\mathbf{x}_i - \mathbf{x}_k\|^2 / 2\sigma_i^2\big)} ,
  \label{eq:tsne_p}                      % neighbour probabilities in the latent space
\end{equation}
where the width $\sigma_i$ is chosen so that the distribution has a given \emph{perplexity},
roughly the effective number of neighbours. In the projection, similarities $q_{ij}$ follow a
heavy-tailed Student-$t$ distribution, $q_{ij} \propto (1 + \|\mathbf{y}_i - \mathbf{y}_j\|^2)^{-1}$,
and the points $\mathbf{y}_i$ are moved by gradient descent to minimize the Kullback--Leibler
divergence between the symmetrized $p_{ij}$ and $q_{ij}$. The heavy tail lets dissimilar trees
drift apart, which gives well-separated groups, but the sizes of the groups and the distances
between them are not meaningful. We ran t-SNE with PCA initialization and 1000 iterations, for
perplexities from 5 to 50 in steps of 5.

\paragraph{UMAP}
Uniform Manifold Approximation and Projection (UMAP) \cite{mcinnes2018umap} is a nonlinear
dimensionality reduction method that seeks to preserve the neighbourhood structure of the
high-dimensional data in a lower-dimensional representation. Given the latent codes
$\mathbf{x}_1,\ldots,\mathbf{x}_N$, UMAP first identifies the nearest neighbours of each
observation using the Euclidean distance. These neighbourhood relationships are represented by
a weighted graph.

For each tree $i$, UMAP determines a local connectivity distance $\rho_i$, corresponding to the
distance to the closest neighbour, and a local scale parameter $\sigma_i$. The membership strength
of the connection between tree $i$ and one of its neighbours $j$ is defined as
\begin{equation}
    w_{ij} =
    \begin{cases}
        1, & d(\mathbf{x}_i,\mathbf{x}_j) \leq \rho_i,\\[4pt]
        \exp\!\left(
        -\dfrac{d(\mathbf{x}_i,\mathbf{x}_j)-\rho_i}{\sigma_i}
        \right),
        & d(\mathbf{x}_i,\mathbf{x}_j) > \rho_i,
    \end{cases}
    \label{eq:umap_weight}
\end{equation}
where $d(\mathbf{x}_i,\mathbf{x}_j)$ denotes the Euclidean distance between two latent codes.
Thus, nearby observations receive a high connection strength, while more distant neighbours have
smaller weights. The local scales are chosen so that the graph adapts to the density of the data
around each observation.

UMAP then seeks a low-dimensional representation
$\mathbf{y}_1,\ldots,\mathbf{y}_N$ whose neighbourhood structure reproduces that of the original
graph. In the low-dimensional space, the strength of an association between two observations is
modelled by
\begin{equation}
    \tilde{w}_{ij}
    =
    \frac{1}
    {1+a\|\mathbf{y}_i-\mathbf{y}_j\|_2^{2b}},
    \label{eq:umap_lowdim}
\end{equation}
where $a$ and $b$ determine the shape of the low-dimensional decay function. The embedding is
obtained by minimizing the cross-entropy between the high- and low-dimensional fuzzy graphs:
\begin{equation}
    \mathcal{L}_{\mathrm{UMAP}}
    =
    -\sum_{i,j}
    \left[
        w_{ij}\log(\tilde{w}_{ij})
        +
        (1-w_{ij})\log(1-\tilde{w}_{ij})
    \right].
    \label{eq:umap_loss}
\end{equation}
Minimizing this objective attracts observations that are strongly connected in the original
latent space and repels observations that have weak or no connections. The resulting embedding
therefore preserves neighbourhood relationships rather than the exact pairwise distances of the
original space.

The scale of the neighbourhood structure is controlled by the number of neighbours considered for
each observation. Small values emphasize local structure, whereas larger values allow broader
relationships to influence the embedding. A minimum-distance parameter further controls how
closely neighbouring observations can be placed in the low-dimensional space.

In our experiments, UMAP was applied directly to the 256-dimensional latent representations using
Euclidean distance and the default minimum-distance value. The neighbourhood size was varied
from 5 to 55 in increments of 5, and both two- and three-dimensional embeddings were generated.

% \paragraph{Parameter sweeps}             % trust what is stable across settings
% Because t-SNE and UMAP depend on their parameters and on the random initialization (fixed with a
% seed), we computed them over the whole range of perplexities and neighbour counts, in 2D and 3D.
% A grouping that appears for every setting reflects the data; one that appears for a single
% setting is more likely an artefact of the projection.
    % Subsection "Visualizing the latent space"
% Results of the clustering -- placeholders for figures, tables and interpretation.
% Replace each \figplaceholder by \includegraphics and each \hint by your own text.

\subsection{Results}                     % what the clustering found
\label{sec:clust_results}                % \secref{sec:clust_results}

\subsubsection{Agglomerative clustering}  % dendrogram, elbow, projections
\label{sec:clust_results_agglo}

\begin{figure}[htbp]                     % dendrogram + elbow curve
  \centering
  \figplaceholder[5cm]{Dendrogram of the Ward clustering (last 30 merges), with the cut line}
  \par\medskip
  \figplaceholder[4cm]{Average intra-cluster distance $D_{\mathrm{avg}}(k)$ and its relative reduction}
  \caption{Agglomerative clustering of the latent codes. Top: dendrogram, cut at $k^{*}$
           clusters (dashed line). Bottom: average intra-cluster distance as a function of the
           number of clusters (Eq.~\eqref{eq:davg}) and its relative reduction
           (Eq.~\eqref{eq:kstar}).}
  \label{fig:clust_agglo}
\end{figure}

\hint{Interpret: how many clusters are suggested, how clear is the largest drop compared with the
next candidates, and what are the cluster sizes? Does the dendrogram show well-separated
branches, or merges at similar heights?}

\begin{figure}[htbp]                     % the three projections, coloured by Ward cluster
  \centering
  \figplaceholder[4.5cm]{Kernel PCA \quad|\quad t-SNE \quad|\quad UMAP, coloured by Ward cluster}
  \caption{Latent codes projected with kernel PCA, t-SNE (perplexity \TODO{value}) and UMAP
           (\texttt{n\_neighbors} \TODO{value}), coloured by agglomerative cluster.}
  \label{fig:clust_agglo_proj}
\end{figure}

\hint{Interpret: are the clusters separated in all three projections, only in some, or in none?
Is the picture stable over the parameter sweep?}

\subsubsection{K-means and kernel K-means}  % choice of K, agreement between methods
\label{sec:clust_results_kmeans}

\begin{figure}[htbp]                     % elbow + silhouette for both methods
  \centering
  \figplaceholder[4cm]{K-means and kernel K-means: inertia (elbow) and mean silhouette for $K = 2, \dots, 10$}
  \caption{Choice of the number of clusters for K-means and kernel K-means: inertia
           (Eq.~\eqref{eq:kmeans}) and mean silhouette (Eq.~\eqref{eq:silhouette}).}
  \label{fig:clust_kmeans_k}
\end{figure}

\begin{table}[htbp]                      % agreement between the partitions
  \centering
  \small
  \caption{Agreement between partitions, measured with the adjusted Rand index (1 = same
           partition, 0 = chance level).}
  \label{tab:clust_ari}
  \begin{tabular}{llc}
    \toprule
    Comparison                                   & Latent codes          & ARI \\
    \midrule
    Agglomerative ($K=4$) vs K-means ($K=3$)     & unregistered          & \TODO{?} \\
    Agglomerative ($K=4$) vs K-means ($K=3$)     & registered            & \TODO{?} \\
    K-means: unregistered vs registered          & --                    & \TODO{?} \\
    Agglomerative: unregistered vs registered    & --                    & \TODO{?} \\
    \bottomrule
  \end{tabular}
\end{table}

\hint{Interpret: which $K$ do the elbow and the silhouette suggest, and do they agree? How much do
the two methods, and the two sets of codes, agree (Table~\ref{tab:clust_ari})? Report the robustness checks of
\secref{sec:clust_kmeans}: seed stability (mean ARI), sensitivity to $\gamma$, and the silhouette
of the real data compared with the no-cluster Gaussian baseline.}

\subsubsection{Morphological characterization of the clusters}  % what distinguishes the groups
\label{sec:clust_results_morph}

For every tree, the six descriptors of \secref{sec:morphometrics} are averaged over its branches
and complemented with two whole-tree counts, the number of branches and the deepest generation.
Comparing these eight values between clusters shows what distinguishes the groups, and colouring
the projections by each value shows whether the latent space is organized along these
anatomical properties.

\begin{figure}[htbp]                     % projections coloured by metric
  \centering
  \figplaceholder[5cm]{Projections coloured by each of the 8 metrics}
  \caption{UMAP projection of the latent codes coloured by the mean value of each
           morphological descriptor and by the two tree counts.}
  \label{fig:clust_metric_proj}
\end{figure}

\begin{figure}[htbp]                     % box plots per cluster
  \centering
  \figplaceholder[5cm]{Box plots of the 8 metrics per cluster}
  \caption{Distribution of the eight metrics in each cluster.}
  \label{fig:clust_boxplots}
\end{figure}

Tables~\ref{tab:unreg-kmeans} to~\ref{tab:reg-agglo} give, for each method and each set of
latent codes, the distribution of the eight metrics within every cluster: mean, standard deviation
and variance, median and interquartile range. The shading of each cell shows the standardized
difference between the cluster mean and the mean over all 419 trees,
\begin{equation}
  d = \frac{\bar{x}_{\text{cluster}} - \bar{x}_{\text{all}}}{s_{\text{all}}} ,
  \label{eq:cluster_d}                   % effect size used for the shading
\end{equation}
where $s_{\text{all}}$ is the standard deviation over all trees: red cells lie above the overall
mean, blue cells below, and darker cells further from it. Expressing the difference in units of
the overall spread makes metrics with very different scales comparable, from tortuosity values
around 0.6 to branch counts above 200. Whether a metric differs significantly between the
clusters can be tested with a Kruskal--Wallis test \cite{kruskal1952}, which compares the
distributions without assuming normality. \TODO{Add the Kruskal--Wallis $p$-values per metric,
e.g.\ as one sentence per table or as an extra column.}

% Per-cluster statistics of the 8 morphological metrics (tables from the clustering notebooks).
% \clustertablesetup (setup/helpers.tex) applies the table spacing locally; \topgap and \botgap
% (same file) are invisible struts that add space above the mean and below the SD line.

\begin{table}[htbp]
\centering
\clustertablesetup
\caption{Morphological metrics per cluster: K-means ($K=3$), latent codes of the unregistered trees.}
\label{tab:unreg-kmeans}
\resizebox{\linewidth}{!}{%
\begin{tabular}{lrrr|r}
\toprule
\topgap\textbf{Metric} & \textcolor{dotA}{\textbullet}\,\textbf{Cluster 0 (n=140)} & \textcolor{dotB}{\textbullet}\,\textbf{Cluster 1 (n=107)} & \textcolor{dotC}{\textbullet}\,\textbf{Cluster 2 (n=172)} & \textbf{All}\botgap \\
\midrule
\topgap\textbf{Divergence} & \cellcolor{loBlue!2}\topgap\textbf{0.3799} & \cellcolor{loBlue!14}\topgap\textbf{0.3645} & \cellcolor{hiRed!10}\topgap\textbf{0.3955} & \topgap\textbf{0.3824} \\
 & \cellcolor{loBlue!2}\stat SD 0.05221 $\cdot$ $\sigma^2$ 0.002726\botgap & \cellcolor{loBlue!14}\stat SD 0.05863 $\cdot$ $\sigma^2$ 0.003437\botgap & \cellcolor{hiRed!10}\stat SD 0.0451 $\cdot$ $\sigma^2$ 0.002034\botgap & \stat SD 0.05258 $\cdot$ $\sigma^2$ 0.002764\botgap \\
\arrayrulecolor{black!20}\hline\arrayrulecolor{black}
\topgap\textbf{Length} & \cellcolor{loBlue!3}\topgap\textbf{52.02} & \cellcolor{loBlue!19}\topgap\textbf{48.82} & \cellcolor{hiRed!14}\topgap\textbf{55.39} & \topgap\textbf{52.58} \\
 & \cellcolor{loBlue!3}\stat SD 7.471 $\cdot$ $\sigma^2$ 55.82\botgap & \cellcolor{loBlue!19}\stat SD 9.631 $\cdot$ $\sigma^2$ 92.77\botgap & \cellcolor{hiRed!14}\stat SD 6.287 $\cdot$ $\sigma^2$ 39.53\botgap & \stat SD 8.077 $\cdot$ $\sigma^2$ 65.23\botgap \\
\arrayrulecolor{black!20}\hline\arrayrulecolor{black}
\topgap\textbf{Complexity} & \cellcolor{loBlue!6}\topgap\textbf{0.7748} & \cellcolor{loBlue!11}\topgap\textbf{0.7643} & \cellcolor{hiRed!12}\topgap\textbf{0.811} & \topgap\textbf{0.787} \\
 & \cellcolor{loBlue!6}\stat SD 0.08176 $\cdot$ $\sigma^2$ 0.006685\botgap & \cellcolor{loBlue!11}\stat SD 0.1027 $\cdot$ $\sigma^2$ 0.01055\botgap & \cellcolor{hiRed!12}\stat SD 0.06113 $\cdot$ $\sigma^2$ 0.003737\botgap & \stat SD 0.08274 $\cdot$ $\sigma^2$ 0.006845\botgap \\
\arrayrulecolor{black!20}\hline\arrayrulecolor{black}
\topgap\textbf{Tortuosity} & \cellcolor{hiRed!11}\topgap\textbf{0.6253} & \cellcolor{loBlue!31}\topgap\textbf{0.6029} & \cellcolor{hiRed!10}\topgap\textbf{0.6247} & \topgap\textbf{0.6193} \\
 & \cellcolor{hiRed!11}\stat SD 0.01658 $\cdot$ $\sigma^2$ $2.75{\times}10^{-4}$\botgap & \cellcolor{loBlue!31}\stat SD 0.02305 $\cdot$ $\sigma^2$ $5.31{\times}10^{-4}$\botgap & \cellcolor{hiRed!10}\stat SD 0.01793 $\cdot$ $\sigma^2$ $3.21{\times}10^{-4}$\botgap & \stat SD 0.02123 $\cdot$ $\sigma^2$ $4.51{\times}10^{-4}$\botgap \\
\arrayrulecolor{black!20}\hline\arrayrulecolor{black}
\topgap\textbf{Stenosis} & \cellcolor{loBlue!7}\topgap\textbf{0.3748} & \cellcolor{hiRed!10}\topgap\textbf{0.3882} & \cellcolor{loBlue!1}\topgap\textbf{0.3799} & \topgap\textbf{0.3803} \\
 & \cellcolor{loBlue!7}\stat SD 0.03389 $\cdot$ $\sigma^2$ 0.001149\botgap & \cellcolor{hiRed!10}\stat SD 0.03415 $\cdot$ $\sigma^2$ 0.001166\botgap & \cellcolor{loBlue!1}\stat SD 0.0248 $\cdot$ $\sigma^2$ $6.15{\times}10^{-4}$\botgap & \stat SD 0.03091 $\cdot$ $\sigma^2$ $9.55{\times}10^{-4}$\botgap \\
\arrayrulecolor{black!20}\hline\arrayrulecolor{black}
\topgap\textbf{Ectasia} & \cellcolor{hiRed!1}\topgap\textbf{0.754} & \cellcolor{loBlue!28}\topgap\textbf{0.7433} & \cellcolor{hiRed!16}\topgap\textbf{0.7597} & \topgap\textbf{0.7536} \\
 & \cellcolor{hiRed!1}\stat SD 0.01499 $\cdot$ $\sigma^2$ $2.25{\times}10^{-4}$\botgap & \cellcolor{loBlue!28}\stat SD 0.01407 $\cdot$ $\sigma^2$ $1.98{\times}10^{-4}$\botgap & \cellcolor{hiRed!16}\stat SD 0.01118 $\cdot$ $\sigma^2$ $1.25{\times}10^{-4}$\botgap & \stat SD 0.01478 $\cdot$ $\sigma^2$ $2.18{\times}10^{-4}$\botgap \\
\arrayrulecolor{black!20}\hline\arrayrulecolor{black}
\topgap\textbf{Number of branches} & \cellcolor{loBlue!9}\topgap\textbf{217} & \topgap\textbf{234.6} & \cellcolor{hiRed!7}\topgap\textbf{250} & \topgap\textbf{235} \\
 & \cellcolor{loBlue!9}\stat SD 78.65 $\cdot$ $\sigma^2$ 6186\botgap & \stat SD 101.7 $\cdot$ $\sigma^2$ 10330\botgap & \cellcolor{hiRed!7}\stat SD 69.41 $\cdot$ $\sigma^2$ 4817\botgap & \stat SD 82.76 $\cdot$ $\sigma^2$ 6850\botgap \\
\arrayrulecolor{black!20}\hline\arrayrulecolor{black}
\topgap\textbf{Deepest generation} & \cellcolor{loBlue!6}\topgap\textbf{10.99} & \cellcolor{hiRed!2}\topgap\textbf{11.3} & \cellcolor{hiRed!4}\topgap\textbf{11.36} & \topgap\textbf{11.22} \\
 & \cellcolor{loBlue!6}\stat SD 1.476 $\cdot$ $\sigma^2$ 2.18\botgap & \cellcolor{hiRed!2}\stat SD 1.829 $\cdot$ $\sigma^2$ 3.344\botgap & \cellcolor{hiRed!4}\stat SD 1.279 $\cdot$ $\sigma^2$ 1.635\botgap & \stat SD 1.506 $\cdot$ $\sigma^2$ 2.269\botgap \\
\bottomrule
\end{tabular}}
\par\smallskip\footnotesize Each cell: mean (bold), standard deviation and variance. Shading: standardized difference $d=(\bar{x}_{\text{cluster}}-\bar{x}_{\text{all}})/s_{\text{all}}$, clipped at $|d|=1$; red = cluster mean above the overall mean, blue = below, darker = larger $|d|$.
\end{table}

\begin{table}[htbp]
\centering
\clustertablesetup
\caption{Morphological metrics per cluster: agglomerative clustering ($K=4$), latent codes of the unregistered trees.}
\label{tab:unreg-agglo}
\resizebox{\linewidth}{!}{%
\begin{tabular}{lrrrr|r}
\toprule
\topgap\textbf{Metric} & \textcolor{dotA}{\textbullet}\,\textbf{Cluster 0 (n=187)} & \textcolor{dotB}{\textbullet}\,\textbf{Cluster 1 (n=112)} & \textcolor{dotC}{\textbullet}\,\textbf{Cluster 2 (n=68)} & \textcolor{dotD}{\textbullet}\,\textbf{Cluster 3 (n=52)} & \textbf{All}\botgap \\
\midrule
\topgap\textbf{Divergence} & \topgap\textbf{0.3824} & \cellcolor{hiRed!15}\topgap\textbf{0.4018} & \cellcolor{loBlue!19}\topgap\textbf{0.3579} & \cellcolor{loBlue!8}\topgap\textbf{0.3724} & \topgap\textbf{0.3824} \\
 & \stat SD 0.05107 $\cdot$ $\sigma^2$ 0.002608\botgap & \cellcolor{hiRed!15}\stat SD 0.03565 $\cdot$ $\sigma^2$ 0.001271\botgap & \cellcolor{loBlue!19}\stat SD 0.06208 $\cdot$ $\sigma^2$ 0.003854\botgap & \cellcolor{loBlue!8}\stat SD 0.05986 $\cdot$ $\sigma^2$ 0.003583\botgap & \stat SD 0.05258 $\cdot$ $\sigma^2$ 0.002764\botgap \\
\arrayrulecolor{black!20}\hline\arrayrulecolor{black}
\topgap\textbf{Length} & \cellcolor{hiRed!1}\topgap\textbf{52.82} & \cellcolor{hiRed!18}\topgap\textbf{56.15} & \cellcolor{loBlue!21}\topgap\textbf{48.37} & \cellcolor{loBlue!15}\topgap\textbf{49.58} & \topgap\textbf{52.58} \\
 & \cellcolor{hiRed!1}\stat SD 6.997 $\cdot$ $\sigma^2$ 48.95\botgap & \cellcolor{hiRed!18}\stat SD 5.685 $\cdot$ $\sigma^2$ 32.32\botgap & \cellcolor{loBlue!21}\stat SD 10.15 $\cdot$ $\sigma^2$ 103.1\botgap & \cellcolor{loBlue!15}\stat SD 9.452 $\cdot$ $\sigma^2$ 89.34\botgap & \stat SD 8.077 $\cdot$ $\sigma^2$ 65.23\botgap \\
\arrayrulecolor{black!20}\hline\arrayrulecolor{black}
\topgap\textbf{Complexity} & \cellcolor{hiRed!1}\topgap\textbf{0.79} & \cellcolor{hiRed!15}\topgap\textbf{0.8175} & \cellcolor{loBlue!16}\topgap\textbf{0.7535} & \cellcolor{loBlue!16}\topgap\textbf{0.754} & \topgap\textbf{0.787} \\
 & \cellcolor{hiRed!1}\stat SD 0.07292 $\cdot$ $\sigma^2$ 0.005317\botgap & \cellcolor{hiRed!15}\stat SD 0.05175 $\cdot$ $\sigma^2$ 0.002678\botgap & \cellcolor{loBlue!16}\stat SD 0.1084 $\cdot$ $\sigma^2$ 0.01175\botgap & \cellcolor{loBlue!16}\stat SD 0.1037 $\cdot$ $\sigma^2$ 0.01076\botgap & \stat SD 0.08274 $\cdot$ $\sigma^2$ 0.006845\botgap \\
\arrayrulecolor{black!20}\hline\arrayrulecolor{black}
\topgap\textbf{Tortuosity} & \cellcolor{loBlue!2}\topgap\textbf{0.6185} & \cellcolor{hiRed!14}\topgap\textbf{0.6265} & \cellcolor{loBlue!23}\topgap\textbf{0.607} & \cellcolor{hiRed!7}\topgap\textbf{0.623} & \topgap\textbf{0.6193} \\
 & \cellcolor{loBlue!2}\stat SD 0.02183 $\cdot$ $\sigma^2$ $4.76{\times}10^{-4}$\botgap & \cellcolor{hiRed!14}\stat SD 0.01584 $\cdot$ $\sigma^2$ $2.51{\times}10^{-4}$\botgap & \cellcolor{loBlue!23}\stat SD 0.02483 $\cdot$ $\sigma^2$ $6.17{\times}10^{-4}$\botgap & \cellcolor{hiRed!7}\stat SD 0.01636 $\cdot$ $\sigma^2$ $2.68{\times}10^{-4}$\botgap & \stat SD 0.02123 $\cdot$ $\sigma^2$ $4.51{\times}10^{-4}$\botgap \\
\arrayrulecolor{black!20}\hline\arrayrulecolor{black}
\topgap\textbf{Stenosis} & \cellcolor{hiRed!2}\topgap\textbf{0.3816} & \cellcolor{loBlue!1}\topgap\textbf{0.3796} & \cellcolor{hiRed!3}\topgap\textbf{0.3829} & \cellcolor{loBlue!8}\topgap\textbf{0.3739} & \topgap\textbf{0.3803} \\
 & \cellcolor{hiRed!2}\stat SD 0.0283 $\cdot$ $\sigma^2$ $8.01{\times}10^{-4}$\botgap & \cellcolor{loBlue!1}\stat SD 0.02333 $\cdot$ $\sigma^2$ $5.44{\times}10^{-4}$\botgap & \cellcolor{hiRed!3}\stat SD 0.04191 $\cdot$ $\sigma^2$ 0.001756\botgap & \cellcolor{loBlue!8}\stat SD 0.03687 $\cdot$ $\sigma^2$ 0.001359\botgap & \stat SD 0.03091 $\cdot$ $\sigma^2$ $9.55{\times}10^{-4}$\botgap \\
\arrayrulecolor{black!20}\hline\arrayrulecolor{black}
\topgap\textbf{Ectasia} & \cellcolor{loBlue!4}\topgap\textbf{0.7522} & \cellcolor{hiRed!19}\topgap\textbf{0.7605} & \cellcolor{loBlue!15}\topgap\textbf{0.7481} & \cellcolor{loBlue!7}\topgap\textbf{0.7509} & \topgap\textbf{0.7536} \\
 & \cellcolor{loBlue!4}\stat SD 0.01548 $\cdot$ $\sigma^2$ $2.4{\times}10^{-4}$\botgap & \cellcolor{hiRed!19}\stat SD 0.009722 $\cdot$ $\sigma^2$ $9.45{\times}10^{-5}$\botgap & \cellcolor{loBlue!15}\stat SD 0.01402 $\cdot$ $\sigma^2$ $1.96{\times}10^{-4}$\botgap & \cellcolor{loBlue!7}\stat SD 0.01706 $\cdot$ $\sigma^2$ $2.91{\times}10^{-4}$\botgap & \stat SD 0.01478 $\cdot$ $\sigma^2$ $2.18{\times}10^{-4}$\botgap \\
\arrayrulecolor{black!20}\hline\arrayrulecolor{black}
\topgap\textbf{Number of branches} & \topgap\textbf{236} & \cellcolor{hiRed!9}\topgap\textbf{253.1} & \cellcolor{loBlue!6}\topgap\textbf{222.9} & \cellcolor{loBlue!13}\topgap\textbf{208.4} & \topgap\textbf{235} \\
 & \stat SD 81.53 $\cdot$ $\sigma^2$ 6648\botgap & \cellcolor{hiRed!9}\stat SD 58.7 $\cdot$ $\sigma^2$ 3445\botgap & \cellcolor{loBlue!6}\stat SD 104.7 $\cdot$ $\sigma^2$ 10960\botgap & \cellcolor{loBlue!13}\stat SD 91.41 $\cdot$ $\sigma^2$ 8356\botgap & \stat SD 82.76 $\cdot$ $\sigma^2$ 6850\botgap \\
\arrayrulecolor{black!20}\hline\arrayrulecolor{black}
\topgap\textbf{Deepest generation} & \cellcolor{hiRed!2}\topgap\textbf{11.29} & \cellcolor{hiRed!6}\topgap\textbf{11.44} & \cellcolor{hiRed!1}\topgap\textbf{11.25} & \cellcolor{loBlue!20}\topgap\textbf{10.48} & \topgap\textbf{11.22} \\
 & \cellcolor{hiRed!2}\stat SD 1.489 $\cdot$ $\sigma^2$ 2.217\botgap & \cellcolor{hiRed!6}\stat SD 1.129 $\cdot$ $\sigma^2$ 1.275\botgap & \cellcolor{hiRed!1}\stat SD 1.872 $\cdot$ $\sigma^2$ 3.504\botgap & \cellcolor{loBlue!20}\stat SD 1.565 $\cdot$ $\sigma^2$ 2.451\botgap & \stat SD 1.506 $\cdot$ $\sigma^2$ 2.269\botgap \\
\bottomrule
\end{tabular}}
\par\smallskip\footnotesize Each cell: mean (bold), standard deviation and variance. Shading: standardized difference $d=(\bar{x}_{\text{cluster}}-\bar{x}_{\text{all}})/s_{\text{all}}$, clipped at $|d|=1$; red = cluster mean above the overall mean, blue = below, darker = larger $|d|$.
\end{table}

\begin{table}[htbp]
\centering
\clustertablesetup
\caption{Morphological metrics per cluster: K-means ($K=3$), latent codes of the registered trees.}
\label{tab:reg-kmeans}
\resizebox{\linewidth}{!}{%
\begin{tabular}{lrrr|r}
\toprule
\topgap\textbf{Metric} & \textcolor{dotA}{\textbullet}\,\textbf{Cluster 0 (n=151)} & \textcolor{dotB}{\textbullet}\,\textbf{Cluster 1 (n=200)} & \textcolor{dotC}{\textbullet}\,\textbf{Cluster 2 (n=68)} & \textbf{All}\botgap \\
\midrule
\topgap\textbf{Divergence} & \cellcolor{hiRed!16}\topgap\textbf{0.4032} & \cellcolor{loBlue!11}\topgap\textbf{0.3676} & \cellcolor{loBlue!2}\topgap\textbf{0.3794} & \topgap\textbf{0.3824} \\
 & \cellcolor{hiRed!16}\stat SD 0.0438 $\cdot$ $\sigma^2$ 0.001918\botgap & \cellcolor{loBlue!11}\stat SD 0.0522 $\cdot$ $\sigma^2$ 0.002725\botgap & \cellcolor{loBlue!2}\stat SD 0.05677 $\cdot$ $\sigma^2$ 0.003222\botgap & \stat SD 0.05258 $\cdot$ $\sigma^2$ 0.002764\botgap \\
\arrayrulecolor{black!20}\hline\arrayrulecolor{black}
\topgap\textbf{Length} & \cellcolor{hiRed!18}\topgap\textbf{56.19} & \cellcolor{loBlue!14}\topgap\textbf{49.78} & \cellcolor{hiRed!1}\topgap\textbf{52.85} & \topgap\textbf{52.58} \\
 & \cellcolor{hiRed!18}\stat SD 6.06 $\cdot$ $\sigma^2$ 36.73\botgap & \cellcolor{loBlue!14}\stat SD 8.307 $\cdot$ $\sigma^2$ 69\botgap & \cellcolor{hiRed!1}\stat SD 8.182 $\cdot$ $\sigma^2$ 66.94\botgap & \stat SD 8.077 $\cdot$ $\sigma^2$ 65.23\botgap \\
\arrayrulecolor{black!20}\hline\arrayrulecolor{black}
\topgap\textbf{Complexity} & \cellcolor{hiRed!16}\topgap\textbf{0.821} & \cellcolor{loBlue!12}\topgap\textbf{0.7613} & \topgap\textbf{0.7869} & \topgap\textbf{0.787} \\
 & \cellcolor{hiRed!16}\stat SD 0.06085 $\cdot$ $\sigma^2$ 0.003703\botgap & \cellcolor{loBlue!12}\stat SD 0.08631 $\cdot$ $\sigma^2$ 0.00745\botgap & \stat SD 0.08787 $\cdot$ $\sigma^2$ 0.007722\botgap & \stat SD 0.08274 $\cdot$ $\sigma^2$ 0.006845\botgap \\
\arrayrulecolor{black!20}\hline\arrayrulecolor{black}
\topgap\textbf{Tortuosity} & \cellcolor{hiRed!5}\topgap\textbf{0.622} & \cellcolor{hiRed!1}\topgap\textbf{0.6197} & \cellcolor{loBlue!13}\topgap\textbf{0.6124} & \topgap\textbf{0.6193} \\
 & \cellcolor{hiRed!5}\stat SD 0.01903 $\cdot$ $\sigma^2$ $3.62{\times}10^{-4}$\botgap & \cellcolor{hiRed!1}\stat SD 0.02246 $\cdot$ $\sigma^2$ $5.04{\times}10^{-4}$\botgap & \cellcolor{loBlue!13}\stat SD 0.02094 $\cdot$ $\sigma^2$ $4.39{\times}10^{-4}$\botgap & \stat SD 0.02123 $\cdot$ $\sigma^2$ $4.51{\times}10^{-4}$\botgap \\
\arrayrulecolor{black!20}\hline\arrayrulecolor{black}
\topgap\textbf{Stenosis} & \topgap\textbf{0.38} & \cellcolor{loBlue!4}\topgap\textbf{0.3772} & \cellcolor{hiRed!13}\topgap\textbf{0.3902} & \topgap\textbf{0.3803} \\
 & \stat SD 0.02613 $\cdot$ $\sigma^2$ $6.83{\times}10^{-4}$\botgap & \cellcolor{loBlue!4}\stat SD 0.03422 $\cdot$ $\sigma^2$ 0.001171\botgap & \cellcolor{hiRed!13}\stat SD 0.02871 $\cdot$ $\sigma^2$ $8.24{\times}10^{-4}$\botgap & \stat SD 0.03091 $\cdot$ $\sigma^2$ $9.55{\times}10^{-4}$\botgap \\
\arrayrulecolor{black!20}\hline\arrayrulecolor{black}
\topgap\textbf{Ectasia} & \cellcolor{hiRed!14}\topgap\textbf{0.7589} & \cellcolor{loBlue!7}\topgap\textbf{0.751} & \cellcolor{loBlue!11}\topgap\textbf{0.7494} & \topgap\textbf{0.7536} \\
 & \cellcolor{hiRed!14}\stat SD 0.01328 $\cdot$ $\sigma^2$ $1.76{\times}10^{-4}$\botgap & \cellcolor{loBlue!7}\stat SD 0.01537 $\cdot$ $\sigma^2$ $2.36{\times}10^{-4}$\botgap & \cellcolor{loBlue!11}\stat SD 0.01278 $\cdot$ $\sigma^2$ $1.63{\times}10^{-4}$\botgap & \stat SD 0.01478 $\cdot$ $\sigma^2$ $2.18{\times}10^{-4}$\botgap \\
\arrayrulecolor{black!20}\hline\arrayrulecolor{black}
\topgap\textbf{Number of branches} & \cellcolor{hiRed!14}\topgap\textbf{263.8} & \cellcolor{loBlue!11}\topgap\textbf{212.9} & \cellcolor{hiRed!1}\topgap\textbf{236.4} & \topgap\textbf{235} \\
 & \cellcolor{hiRed!14}\stat SD 74.84 $\cdot$ $\sigma^2$ 5601\botgap & \cellcolor{loBlue!11}\stat SD 77.67 $\cdot$ $\sigma^2$ 6032\botgap & \cellcolor{hiRed!1}\stat SD 94.54 $\cdot$ $\sigma^2$ 8938\botgap & \stat SD 82.76 $\cdot$ $\sigma^2$ 6850\botgap \\
\arrayrulecolor{black!20}\hline\arrayrulecolor{black}
\topgap\textbf{Deepest generation} & \cellcolor{hiRed!8}\topgap\textbf{11.53} & \cellcolor{loBlue!7}\topgap\textbf{10.96} & \cellcolor{hiRed!2}\topgap\textbf{11.31} & \topgap\textbf{11.22} \\
 & \cellcolor{hiRed!8}\stat SD 1.356 $\cdot$ $\sigma^2$ 1.837\botgap & \cellcolor{loBlue!7}\stat SD 1.49 $\cdot$ $\sigma^2$ 2.219\botgap & \cellcolor{hiRed!2}\stat SD 1.739 $\cdot$ $\sigma^2$ 3.023\botgap & \stat SD 1.506 $\cdot$ $\sigma^2$ 2.269\botgap \\
\bottomrule
\end{tabular}}
\par\smallskip\footnotesize Each cell: mean (bold), standard deviation and variance. Shading: standardized difference $d=(\bar{x}_{\text{cluster}}-\bar{x}_{\text{all}})/s_{\text{all}}$, clipped at $|d|=1$; red = cluster mean above the overall mean, blue = below, darker = larger $|d|$.
\end{table}

\begin{table}[htbp]
\centering
\clustertablesetup
\caption{Morphological metrics per cluster: agglomerative clustering ($K=4$), latent codes of the registered trees.}
\label{tab:reg-agglo}
\resizebox{\linewidth}{!}{%
\begin{tabular}{lrrrr|r}
\toprule
\topgap\textbf{Metric} & \textcolor{dotA}{\textbullet}\,\textbf{Cluster 0 (n=325)} & \textcolor{dotB}{\textbullet}\,\textbf{Cluster 1 (n=74)} & \textcolor{dotC}{\textbullet}\,\textbf{Cluster 2 (n=13)} & \textcolor{dotD}{\textbullet}\,\textbf{Cluster 3 (n=7)} & \textbf{All}\botgap \\
\midrule
\topgap\textbf{Divergence} & \cellcolor{loBlue!3}\topgap\textbf{0.3789} & \cellcolor{hiRed!8}\topgap\textbf{0.3931} & \cellcolor{hiRed!2}\topgap\textbf{0.3854} & \cellcolor{hiRed!32}\topgap\textbf{0.4242} & \topgap\textbf{0.3824} \\
 & \cellcolor{loBlue!3}\stat SD 0.05366 $\cdot$ $\sigma^2$ 0.002879\botgap & \cellcolor{hiRed!8}\stat SD 0.04695 $\cdot$ $\sigma^2$ 0.002205\botgap & \cellcolor{hiRed!2}\stat SD 0.04919 $\cdot$ $\sigma^2$ 0.00242\botgap & \cellcolor{hiRed!32}\stat SD 0.03629 $\cdot$ $\sigma^2$ 0.001317\botgap & \stat SD 0.05258 $\cdot$ $\sigma^2$ 0.002764\botgap \\
\arrayrulecolor{black!20}\hline\arrayrulecolor{black}
\topgap\textbf{Length} & \cellcolor{loBlue!2}\topgap\textbf{52.27} & \cellcolor{hiRed!4}\topgap\textbf{53.34} & \cellcolor{loBlue!1}\topgap\textbf{52.42} & \cellcolor{hiRed!35}\topgap\textbf{59.61} & \topgap\textbf{52.58} \\
 & \cellcolor{loBlue!2}\stat SD 8.174 $\cdot$ $\sigma^2$ 66.82\botgap & \cellcolor{hiRed!4}\stat SD 6.939 $\cdot$ $\sigma^2$ 48.15\botgap & \cellcolor{loBlue!1}\stat SD 11.04 $\cdot$ $\sigma^2$ 121.9\botgap & \cellcolor{hiRed!35}\stat SD 6.027 $\cdot$ $\sigma^2$ 36.33\botgap & \stat SD 8.077 $\cdot$ $\sigma^2$ 65.23\botgap \\
\arrayrulecolor{black!20}\hline\arrayrulecolor{black}
\topgap\textbf{Complexity} & \cellcolor{loBlue!2}\topgap\textbf{0.7822} & \cellcolor{hiRed!8}\topgap\textbf{0.8042} & \cellcolor{loBlue!8}\topgap\textbf{0.7707} & \cellcolor{hiRed!35}\topgap\textbf{0.8591} & \topgap\textbf{0.787} \\
 & \cellcolor{loBlue!2}\stat SD 0.08453 $\cdot$ $\sigma^2$ 0.007145\botgap & \cellcolor{hiRed!8}\stat SD 0.06708 $\cdot$ $\sigma^2$ 0.004499\botgap & \cellcolor{loBlue!8}\stat SD 0.1059 $\cdot$ $\sigma^2$ 0.01121\botgap & \cellcolor{hiRed!35}\stat SD 0.04858 $\cdot$ $\sigma^2$ 0.00236\botgap & \stat SD 0.08274 $\cdot$ $\sigma^2$ 0.006845\botgap \\
\arrayrulecolor{black!20}\hline\arrayrulecolor{black}
\topgap\textbf{Tortuosity} & \topgap\textbf{0.6195} & \cellcolor{loBlue!3}\topgap\textbf{0.6177} & \cellcolor{hiRed!8}\topgap\textbf{0.6235} & \cellcolor{hiRed!3}\topgap\textbf{0.621} & \topgap\textbf{0.6193} \\
 & \stat SD 0.02139 $\cdot$ $\sigma^2$ $4.58{\times}10^{-4}$\botgap & \cellcolor{loBlue!3}\stat SD 0.02139 $\cdot$ $\sigma^2$ $4.57{\times}10^{-4}$\botgap & \cellcolor{hiRed!8}\stat SD 0.0197 $\cdot$ $\sigma^2$ $3.88{\times}10^{-4}$\botgap & \cellcolor{hiRed!3}\stat SD 0.01674 $\cdot$ $\sigma^2$ $2.8{\times}10^{-4}$\botgap & \stat SD 0.02123 $\cdot$ $\sigma^2$ $4.51{\times}10^{-4}$\botgap \\
\arrayrulecolor{black!20}\hline\arrayrulecolor{black}
\topgap\textbf{Stenosis} & \topgap\textbf{0.3802} & \cellcolor{loBlue!1}\topgap\textbf{0.3797} & \cellcolor{loBlue!2}\topgap\textbf{0.3789} & \cellcolor{hiRed!17}\topgap\textbf{0.3937} & \topgap\textbf{0.3803} \\
 & \stat SD 0.03151 $\cdot$ $\sigma^2$ $9.93{\times}10^{-4}$\botgap & \cellcolor{loBlue!1}\stat SD 0.02922 $\cdot$ $\sigma^2$ $8.54{\times}10^{-4}$\botgap & \cellcolor{loBlue!2}\stat SD 0.03077 $\cdot$ $\sigma^2$ $9.47{\times}10^{-4}$\botgap & \cellcolor{hiRed!17}\stat SD 0.02064 $\cdot$ $\sigma^2$ $4.26{\times}10^{-4}$\botgap & \stat SD 0.03091 $\cdot$ $\sigma^2$ $9.55{\times}10^{-4}$\botgap \\
\arrayrulecolor{black!20}\hline\arrayrulecolor{black}
\topgap\textbf{Ectasia} & \cellcolor{loBlue!1}\topgap\textbf{0.7532} & \cellcolor{hiRed!1}\topgap\textbf{0.7541} & \cellcolor{hiRed!4}\topgap\textbf{0.755} & \cellcolor{hiRed!31}\topgap\textbf{0.765} & \topgap\textbf{0.7536} \\
 & \cellcolor{loBlue!1}\stat SD 0.01518 $\cdot$ $\sigma^2$ $2.3{\times}10^{-4}$\botgap & \cellcolor{hiRed!1}\stat SD 0.01261 $\cdot$ $\sigma^2$ $1.59{\times}10^{-4}$\botgap & \cellcolor{hiRed!4}\stat SD 0.0177 $\cdot$ $\sigma^2$ $3.13{\times}10^{-4}$\botgap & \cellcolor{hiRed!31}\stat SD 0.006023 $\cdot$ $\sigma^2$ $3.63{\times}10^{-5}$\botgap & \stat SD 0.01478 $\cdot$ $\sigma^2$ $2.18{\times}10^{-4}$\botgap \\
\arrayrulecolor{black!20}\hline\arrayrulecolor{black}
\topgap\textbf{Number of branches} & \cellcolor{loBlue!2}\topgap\textbf{230.4} & \cellcolor{hiRed!8}\topgap\textbf{252} & \cellcolor{loBlue!7}\topgap\textbf{220.5} & \cellcolor{hiRed!31}\topgap\textbf{299.3} & \topgap\textbf{235} \\
 & \cellcolor{loBlue!2}\stat SD 84.18 $\cdot$ $\sigma^2$ 7087\botgap & \cellcolor{hiRed!8}\stat SD 71.89 $\cdot$ $\sigma^2$ 5167\botgap & \cellcolor{loBlue!7}\stat SD 83.75 $\cdot$ $\sigma^2$ 7013\botgap & \cellcolor{hiRed!31}\stat SD 87.31 $\cdot$ $\sigma^2$ 7624\botgap & \stat SD 82.76 $\cdot$ $\sigma^2$ 6850\botgap \\
\arrayrulecolor{black!20}\hline\arrayrulecolor{black}
\topgap\textbf{Deepest generation} & \cellcolor{loBlue!1}\topgap\textbf{11.19} & \cellcolor{hiRed!4}\topgap\textbf{11.38} & \cellcolor{loBlue!6}\topgap\textbf{11} & \cellcolor{hiRed!9}\topgap\textbf{11.57} & \topgap\textbf{11.22} \\
 & \cellcolor{loBlue!1}\stat SD 1.541 $\cdot$ $\sigma^2$ 2.375\botgap & \cellcolor{hiRed!4}\stat SD 1.372 $\cdot$ $\sigma^2$ 1.882\botgap & \cellcolor{loBlue!6}\stat SD 1.528 $\cdot$ $\sigma^2$ 2.333\botgap & \cellcolor{hiRed!9}\stat SD 1.272 $\cdot$ $\sigma^2$ 1.619\botgap & \stat SD 1.506 $\cdot$ $\sigma^2$ 2.269\botgap \\
\bottomrule
\end{tabular}}
\par\smallskip\footnotesize Each cell: mean (bold), standard deviation and variance. Shading: standardized difference $d=(\bar{x}_{\text{cluster}}-\bar{x}_{\text{all}})/s_{\text{all}}$, clipped at $|d|=1$; red = cluster mean above the overall mean, blue = below, darker = larger $|d|$.
\end{table}  % the four landscape tables

\hint{Interpret: which metrics separate the clusters most (darkest cells), and in which
direction? Can the clusters be described in anatomical terms, e.g.\ longer, more complex and
more branched trees against shorter, simpler ones? Compare unregistered and registered codes:
does registration change which properties drive the clusters? Note also the cluster sizes: with
registered codes, agglomerative clustering produces two very small clusters (13 and 7 trees)
next to one of 325 trees. Do the metric-coloured projections show gradients that follow the
clusters, or that cut across them?}        % Subsection "Results" (placeholders)
% \FloatBarrier                              % all clustering figures before the next section
